% Unit 032 terjemahan Bahasa Indonesia dari chapter3.tex baris 57--162.
% Sumber: Methods of Algebra, Volume 2, commit
% 9a5803ff77dd3257484cb177f851a73770a59dd3 (CC BY 4.0).
% stable-unit-id: o014.aljabr2.chapter3.complexes-over-additive-categories
% Cakupan: Bagian 3.1 lengkap.

% segment-id: o014.li2.chapter3.complexes-over-additive-categories.h001
\section{Kompleks pada Kategori Aditif}\label{sec:additive-cplx}
\hypertarget{o014.aljabr2.chapter3.complexes-over-additive-categories}{ }

% segment-id: o014.li2.chapter3.complexes-over-additive-categories.p001
Definisi~\ref{def:complex} telah memperkenalkan secara ringkas konsep kompleks
beserta notasi yang lazim digunakan. Misalkan $\mathcal{A}$ suatu kategori
aditif dan tinjau kompleks tak terbatas pada kategori tersebut, yakni
kompleks tanpa suku ujung
$X=(X^n,d^n)_{n\in\Z}$; morfisme $d^n$ sering dinotasikan dengan $d_X^n$.
Karena alasan historis, morfisme-morfisme $d_X^n$ ini juga disebut
``diferensial''.
\index{diferensial (differential)}

% segment-id: o014.li2.chapter3.complexes-over-additive-categories.d001
\begin{definisi}\label{def:cplx-cat}
	\index{kompleks}
	\index[sym1]{CA@$\cate{C}(\mathcal{A})$}
	Suatu morfisme dari kompleks $X$ ke kompleks $Y$ didefinisikan sebagai
	data $(f^n)_{n\in\Z}$, yang juga disingkat menjadi $f$, dengan
	$f^n\in\Hom_{\mathcal{A}}(X^n,Y^n)$ memenuhi
% segment-id: o014.li2.chapter3.complexes-over-additive-categories.q001
	\[ d_Y^n f^n=f^{n+1}d_X^n,\quad n\in\Z. \]

	Semua kompleks beserta morfisme di antaranya membentuk kategori
	$\cate{C}(\mathcal{A})$. Morfisme identitas suatu kompleks $X$ ditentukan
	oleh $(\identity_X)^n:=\identity_{X^n}$, sedangkan komposisi morfisme
	didefinisikan melalui $(gf)^n=g^nf^n$. Penjumlahan
	$f+g:=(f^n+g^n)_{n\in\Z}$ menjadikan
	$\Hom_{\cate{C}(\mathcal{A})}(X,Y)$ suatu grup abelian.
\end{definisi}

% segment-id: o014.li2.chapter3.complexes-over-additive-categories.p002
Jika $\mathcal{A}$ merupakan kategori $\Bbbk$-linear dalam pengertian
Definisi~\ref{def:k-linear-cat}, dengan $\Bbbk$ suatu gelanggang komutatif
yang telah ditetapkan, maka
$\Hom_{\cate{C}(\mathcal{A})}(X,Y)$ juga merupakan modul-$\Bbbk$ terhadap
perkalian skalar suku demi suku. Dalam semua hasil pada bagian ini, kata
``aditif'' selalu dapat diganti dengan istilah ``$\Bbbk$-linear'' yang lebih
umum. Demi keringkasan, selanjutnya hanya versi aditif yang dinyatakan.

% segment-id: o014.li2.chapter3.complexes-over-additive-categories.p003
Sifat utama dalam definisi kompleks ialah $d^{n+1}d^n=0$. Sifat ini juga dapat
dinyatakan kembali dengan bahasa objek diferensial bergradasi.%
\footnote{Pembahasan terkait akan dilanjutkan dalam
\S~\sourcecrossref{sec:spectral-sequence}{5.2}.}

% segment-id: o014.li2.chapter3.complexes-over-additive-categories.d002
\begin{definisi}[Objek bergradasi]\label{def:graded-obj}
	\index{objek bergradasi (graded object)}
	\index{objek bergradasi!$\Z^m$; ganda (bigraded)}
	Misalkan $\mathcal{A}$ sembarang kategori. Ingat kembali konstruksi kategori
	produk $\mathcal{A}^{\Z}$:
% segment-id: o014.li2.chapter3.complexes-over-additive-categories.l001
	\begin{compactitem}
% segment-id: o014.li2.chapter3.complexes-over-additive-categories.i001
		\item objeknya berbentuk $X=(X^n)_{n\in\Z}$, dengan
		$X^n\in\Obj(\mathcal{A})$;
% segment-id: o014.li2.chapter3.complexes-over-additive-categories.i002
		\item morfismenya berbentuk
		$f=(f^n:X^n\to Y^n)_{n\in\Z}$ tanpa syarat tambahan, dan komposisi
		dilakukan suku demi suku (atau derajat demi derajat), yakni
		$(gf)^n=g^nf^n$;
% segment-id: o014.li2.chapter3.complexes-over-additive-categories.i003
		\item definisikan automorfisme $T$ pada $\mathcal{A}^{\Z}$ sebagai
		``funktor pergeseran'' berikut:
		$(TX)^n=X^{n+1}$ dan $(Tf)^n=f^{n+1}$.
	\end{compactitem}
	Objek $X=(X^n)_n$ juga disebut \emph{objek bergradasi $\Z$} pada
	$\mathcal{A}$, atau singkatnya \emph{objek bergradasi}.

	Secara lebih umum, objek-objek dari $\mathcal{A}^{\Z^m}$ disebut objek
	bergradasi $\Z^m$ dan ditulis
	$(X^{n_1,\ldots,n_m})_{n_1,\ldots,n_m}$; morfismenya ditulis
	$(f^{n_1,\ldots,n_m})_{n_1,\ldots,n_m}$. Kategori ini dilengkapi keluarga
	funktor pergeseran $T_1,\ldots,T_m$ yang saling berkomutasi. Kasus khusus
	$m=2$ juga disebut \emph{objek bergradasi ganda}.
\end{definisi}

% segment-id: o014.li2.chapter3.complexes-over-additive-categories.p004
Jika $\mathcal{A}$ kategori aditif, maka $\mathcal{A}^{\Z^m}$ juga kategori
aditif: penjumlahan morfisme, jumlah langsung objek, dan semua operasi lainnya
dilakukan derajat demi derajat.

% segment-id: o014.li2.chapter3.complexes-over-additive-categories.d003
\begin{definisi}[Objek diferensial bergradasi]
	\index{objek diferensial bergradasi (differential graded object)}
	Misalkan $\mathcal{A}$ kategori aditif. Suatu objek diferensial bergradasi
	pada $\mathcal{A}$ didefinisikan sebagai data $(X,d)$, dengan
	$X\in\Obj(\mathcal{A}^{\Z})$ dan $d:X\to TX$ memenuhi $(Td)d=0$. Suatu
	morfisme $(X,d_X)\to(Y,d_Y)$ didefinisikan sebagai morfisme
	$f\in\Hom_{\mathcal{A}^{\Z}}(X,Y)$ yang memenuhi $(Tf)d_X=d_Yf$.
\end{definisi}

% segment-id: o014.li2.chapter3.complexes-over-additive-categories.p005
Jika $d$ dalam objek diferensial bergradasi $(X,d)$ dituliskan secara konkret
sebagai keluarga morfisme $d_X^n:X^n\to X^{n+1}$, maka syarat pada $d$ ialah
$d_X^{n+1}d_X^n=0$, sedangkan syarat pada morfisme ialah
$f^{n+1}d_X^n=d_Y^nf^n$. Dengan demikian, kita kembali memperoleh
Definisi~\ref{def:cplx-cat}. Jadi, kategori objek diferensial bergradasi pada
$\mathcal{A}$ isomorfik dengan kategori kompleks. Dalam bagian ini kita akan
beralih di antara kedua sudut pandang tersebut tanpa keterangan lebih lanjut.

% segment-id: o014.li2.chapter3.complexes-over-additive-categories.p006
Langkah pertama dalam mengkaji kompleks ialah memahami berbagai limit dalam
$\cate{C}(\mathcal{A})$. Versi untuk kategori produk $\mathcal{A}^{\Z}$
relatif mudah: limitnya dapat dibentuk pada setiap $n\in\Z$, yakni suku demi
suku atau derajat demi derajat di dalam $\mathcal{A}$. Berikut ini dijelaskan
cara mengangkat konstruksi tersebut ke tingkat $\cate{C}(\mathcal{A})$.
Tuliskan $U:\cate{C}(\mathcal{A})\to\mathcal{A}^{\Z}$ untuk funktor pelupa
yang memetakan objek diferensial bergradasi $(X,d)$ ke objek bergradasi $X$.

% segment-id: o014.li2.chapter3.complexes-over-additive-categories.le001
\begin{lema}[Konstruksi limit suku demi suku]\label{prop:limits-cplx}
	Untuk setiap funktor $\alpha:I\to\cate{C}(\mathcal{A})$, tuliskan
	$\overline{\alpha}:=U\alpha$. Andaikan $\varinjlim\overline{\alpha}$ ada.
	Objek ini dapat dilengkapi secara unik dengan struktur objek
	$\cate{C}(\mathcal{A})$ sedemikian sehingga menjadi
	$\varinjlim\alpha$. Pernyataan yang bersesuaian juga berlaku untuk
	$\varprojlim$.

	Dalam bahasa Definisi~\ref{def:create-limit}, pernyataan ini berarti bahwa
	funktor pelupa $U$ menciptakan $\varinjlim$ dan $\varprojlim$.
\end{lema}
% segment-id: o014.li2.chapter3.complexes-over-additive-categories.pf001
\begin{bukti}
	Misalkan funktor $\alpha:I\to\cate{C}(\mathcal{A})$ memetakan
	$i\in\Obj(I)$ ke objek diferensial bergradasi
	$\alpha(i)=(\overline{\alpha}(i),d_i)$. Jika
	$\varinjlim\overline{\alpha}$ ada, semua $d_i$ menentukan
	$d:\varinjlim\overline{\alpha}\to
	(T\varinjlim\overline{\alpha}\simeq
	\varinjlim T\overline{\alpha})$; di sini automorfisme $T$ secara otomatis
	mempertahankan $\varinjlim$. Morfisme tersebut dicirikan oleh diagram
	komutatif
% segment-id: o014.li2.chapter3.complexes-over-additive-categories.g001
	\[\begin{tikzcd}
		\varinjlim \overline{\alpha} \arrow[r, "d"] & \varinjlim T\overline{\alpha} \arrow[r, "\sim"] & T\varinjlim \overline{\alpha} \\
		\overline{\alpha}(i) \arrow[u] \arrow[r, "{d_i}"'] & T\overline{\alpha}(i) \arrow[u] \arrow[ru] &
	\end{tikzcd} \qquad (i \in \Obj(I)) . \]

	Dari $(Td_i)d_i=0$ diperoleh $(Td)d=0$, sehingga
	$(\varinjlim\overline{\alpha},d)$ merupakan objek
	$\cate{C}(\mathcal{A})$. Diagram di atas juga menunjukkan bahwa inilah
	satu-satunya pilihan yang menjadikan semua pemetaan
	$(\overline{\alpha}(i),d_i)\to
	(\varinjlim\overline{\alpha},d)$ sebagai morfisme.

	Sekarang kita verifikasi sifat universalnya sebagai $\varinjlim\alpha$
	dalam $\cate{C}(\mathcal{A})$. Misalkan diberikan keluarga morfisme yang
	kompatibel $f_i:\alpha(i)\to L$ dalam $\cate{C}(\mathcal{A})$. Setelah
	struktur diferensial dilupakan, kita memperoleh
	$f_i:\overline{\alpha}(i)\to U(L)$, yang secara unik menentukan
	$f:\varinjlim\overline{\alpha}\to U(L)$. Hal utama yang perlu dibuktikan
	ialah bahwa $f$ merupakan morfisme dalam $\cate{C}(\mathcal{A})$. Walaupun
	hal ini sudah dapat diduga, kita tuliskan diagramnya secara eksplisit:
% segment-id: o014.li2.chapter3.complexes-over-additive-categories.g002
	\[\begin{tikzcd}
		\varinjlim \overline{\alpha} \arrow[dd, "f"'] \arrow[rr, "d"] & & \varinjlim T\overline{\alpha} \arrow[r, "\sim"] & T\varinjlim \overline{\alpha} \arrow[dd, "{Tf}"] \\
		& \overline{\alpha}(i) \arrow[lu] \arrow[ld, "{f_i}"] \arrow[r, "d_i"] & T\overline{\alpha}(i) \arrow[ru] \arrow[rd, "{Tf_i}"] \arrow[u] & \\
		U(L) \arrow[rrr, "d_L"'] & & & TU(L)
	\end{tikzcd} \qquad (i \in \Obj(I)). \]
	Komutativitas setiap bagiannya dapat diperiksa langsung dari definisi atau
	konstruksi. Jadi, komposisi $d_Lf$ dan $(Tf)d$ dengan setiap
	$\overline{\alpha}(i)\to\varinjlim\overline{\alpha}$ adalah sama. Oleh karena itu,
	$d_Lf=(Tf)d$, sebagaimana hendak dibuktikan.
\end{bukti}

% segment-id: o014.li2.chapter3.complexes-over-additive-categories.pr001
\begin{proposisi}\label{prop:additive-cat-cplx}
	Kategori $\cate{C}(\mathcal{A})$ merupakan kategori aditif. Jika
	$\mathcal{A}$ lengkap (atau kolengkap), maka $\cate{C}(\mathcal{A})$ juga
	lengkap (atau kolengkap).
\end{proposisi}
% segment-id: o014.li2.chapter3.complexes-over-additive-categories.pf002
\begin{bukti}
	Gunakan Lema~\ref{prop:limits-cplx} untuk mereduksi semua limit yang
	diperlukan menjadi limit derajat demi derajat di $\mathcal{A}$.
\end{bukti}

% segment-id: o014.li2.chapter3.complexes-over-additive-categories.d004
\begin{definisi}[Funktor pergeseran]\label{def:translation-functor}
	\index{funktor pergeseran (translation functor)}
	\index[sym1]{[n]@$[n]$}
	Misalkan $X$ objek $\cate{C}(\mathcal{A})$ dan $n\in\Z$. Definisikan
	kompleks $X[n]$ sebagai berikut:%
	\footnote{Penyisipan $(-1)^n$ dalam definisi $d_{X[n]}$ bermanfaat.
	Versi ini juga isomorfik dengan versi tanpa tanda,
	$X[n]^\circ:=\left(X^{n+k},d_X^{k+n}\right)_k$. Untuk kasus $n=1$,
	definisikan $X[1]\rightiso X[1]^\circ$ dengan pola tanda
	$\cdots,+,-,+,\cdots$.}
% segment-id: o014.li2.chapter3.complexes-over-additive-categories.q002
	\[ \left(X[n]\right)^k:=X^{k+n},\quad
	d_{X[n]}^k:=(-1)^nd_X^{k+n}. \]
	Konstruksi ini memberikan funktor aditif
	$[n]:\cate{C}(\mathcal{A})\to\cate{C}(\mathcal{A})$. Jika
	$f:X\to Y$ merupakan morfisme kompleks, maka $f[n]:X[n]\to Y[n]$
	diberikan oleh $f[n]^k:=f^{n+k}$.
\end{definisi}

% segment-id: o014.li2.chapter3.complexes-over-additive-categories.p007
Jelas bahwa $X[n]$ tetap merupakan kompleks, sehingga funktor tersebut
terdefinisi dengan baik. Sifat-sifat berikut juga langsung terlihat.
% segment-id: o014.li2.chapter3.complexes-over-additive-categories.l002
\begin{itemize}
% segment-id: o014.li2.chapter3.complexes-over-additive-categories.i004
	\item $[0]=\identity_{\cate{C}(\mathcal{A})}$.
% segment-id: o014.li2.chapter3.complexes-over-additive-categories.i005
	\item $[n][m]=[n+m]$. Jadi, $[n]$ merupakan automorfisme dari
	$\cate{C}(\mathcal{A})$, dengan invers $[-n]$.
% segment-id: o014.li2.chapter3.complexes-over-additive-categories.i006
	\item Keluarga $(d_X^n)_{n\in\Z}$ memberikan morfisme kompleks
	$d_X:X\to X[1]$ (Petunjuk: $d_X^{n+1}d_X^n=0$).
\end{itemize}

% segment-id: o014.li2.chapter3.complexes-over-additive-categories.c001
\begin{konvensi}\label{con:concentrated-cplx}
	Setiap objek $S$ dari $\mathcal{A}$ dapat dipandang sebagai kompleks
	$(S^n,d_S^n)_{n\in\Z}$ dengan $S^0:=S$, semua suku lainnya sama dengan
	$0$, dan setiap $d_S^n$ merupakan morfisme nol. Identifikasi ini memberikan
	funktor aditif yang penuh dan setia
	$\mathcal{A}\hookrightarrow\cate{C}(\mathcal{A})$. Secara lebih umum,
	untuk setiap $n\in\Z$, $S[-n]$ adalah kompleks yang terkonsentrasi pada
	derajat $n$: kompleks ini menempatkan $S$ pada suku berderajat $n$ dan
	menempatkan $0$ pada semua suku lainnya.
\end{konvensi}

% segment-id: o014.li2.chapter3.complexes-over-additive-categories.p008
Hasil berikut langsung diperoleh.

% segment-id: o014.li2.chapter3.complexes-over-additive-categories.pr002
\begin{proposisi}\label{prop:CA-functoriality}
	\index[sym1]{CF@$\cate{C}F$}
	Untuk suatu funktor aditif $F:\mathcal{A}\to\mathcal{A}'$, pemetaan
	$(X^n,d_X^n)_n\mapsto(FX^n,Fd_X^n)_n$ memberikan funktor
	$\cate{C}F:\cate{C}(\mathcal{A})\to\cate{C}(\mathcal{A}')$. Funktor ini
	berkomutasi dengan funktor pergeseran:
	$[m]\circ\cate{C}F=\cate{C}F\circ[m]$ untuk setiap $m\in\Z$.
\end{proposisi}

% segment-id: o014.li2.chapter3.complexes-over-additive-categories.r001
\begin{catatan}\label{rem:chain-cplx-translation}
	\index{funktor pergeseran}
	\index[sym1]{[n]}
	Semua hasil dalam bagian ini juga berlaku untuk kompleks rantai
	$(X_\bullet,d_\bullet)$ yang disebutkan dalam
	Catatan~\ref{rem:cochain-vs-chain}. Funktor pergeseran pada kategori kompleks
	rantai didefinisikan sebagai berikut:
% segment-id: o014.li2.chapter3.complexes-over-additive-categories.q003
	\[ X[n]_k:=X_{k+n},\quad d^{X[n]}_k=(-1)^nd^X_{n+k} \]
	Jika peralihan antara kedua konvensi dilakukan melalui $X_n=X^{-n}$ dan
	$d_n=d^{-n}$, maka pergeseran $[1]$ untuk kompleks rantai bersesuaian dengan
	pergeseran $[-1]$ untuk kompleks kokrantai.
\end{catatan}
